\chapter{O que o conjunto sustenta}
\label{cap:estabelece}

Nove resultados se sustentam, oito conclusões foram retiradas por análise
posterior e três perguntas ficam em aberto. Cada uma das conclusões retiradas
foi afirmada, em algum momento, com a mesma confiança das que ficaram.

\section{Magnitude e composição do corte}

O corte é grande e de origem sistêmica. A taxa é de 21,6\% no parque com
apuração, dois terços por razão energética, incidindo sobre toda a janela solar.

A geração de referência do ONS é calibrada. O viés mediano é de 1,02 em 74
usinas, e não se degrada onde o corte é mais intenso. O agregado da
seção~\ref{sec:fig1} permanece válido.

A perda de Sol do Cerrado é de rede, e não de recurso. Nove décimos da queda de
29\% no fator de capacidade não se explicam pela irradiância.

O corte pode ser estimado sem a referência do operador. Irradiância e razão de
desempenho calibrada estimam 33,5\% de perda em 2025, e 34,4\% com o passo de
área acoplado, contra 32,8\% apurados pelo ONS.

\begin{objecao}
A independência dessa estimativa é parcial. Ela dispensa a geração de
referência, e ainda usa a apuração do ONS para selecionar os dias de calibração.
O sinalizador de restrição é compartilhado entre as duas.
\end{objecao}

\section{Observabilidade por satélite}

A geração distribuída é sub-pixel. São 91\% das instalações abaixo de um pixel
Sentinel-2, e a verificação por imagem exige 2 m ou melhor, além de coordenada
por instalação, que não existe no registro.

A construção de usina centralizada é detectável e datável; a operação não é. O
brilho relativo sobe 34\% na construção, com $p = 1{,}2 \times 10^{-5}$, e não
apresenta sinal na operação, com $p = 0{,}09$.

\section{Despacho e alocação}

O piso térmico inflexível está no Sudeste, e não no Nordeste. São 1.739 MWmed
de inflexibilidade pura no primeiro, contra 36 no segundo.

O grau do nó está associado ao corte eólico, com $\rho = +0{,}35$. É o único dos
quatro resultados de rede a passar em bootstrap, partição estrutural e troca de
grafo.

O preditor acompanha a origem do corte. A geração solar tem corte 91\% sistêmico
e responde à distância até a demanda. A eólica tem 24\% de corte local e responde
a aglomeração no nó, com $\rho = +0{,}47$. A dissociação se mantém sob controle
de geografia.

\section{As conclusões retiradas}

As oito afirmações abaixo foram sustentadas por uma análise e retiradas por
outra. Elas permanecem no relatório porque a correção é parte do resultado.

\begin{enumerate}[leftmargin=1.2em,itemsep=0.3em]
\item \textbf{A referência do ONS superestima o corte.} Era artefato do
contrafactual, porque 2023 já era um ano com restrição. A revisão reforça a
conclusão principal: o corte implícito é de 33\%, e não de 26\%.

\item \textbf{A térmica inflexível desloca a geração solar.} Essa térmica se
associa igualmente ao corte a dois mil quilômetros de distância. Ela é marcador
sistêmico, e o desenho não separa as hipóteses.

\item \textbf{O Nordeste é cortado porque o corredor satura.} Nas horas de corte
alto o fluxo para o Sudeste é menor que nas demais.

\item \textbf{A térmica não abre espaço para a solar.} A afirmação não sobrevive
ao controle de carga, porque a carga líquida diária é plana.

\item \textbf{Explicações locais do corte generalizam.} Nenhuma das cinco
covariáveis testadas no par se reproduz no parque.

\item \textbf{A distância à demanda vale para qualquer tecnologia.} A replicação
em eólica restringe o enunciado a fonte de corte sistêmico, e os critérios de
robustez mostram que, mesmo na solar, o resultado é o menos sustentado dos
quatro.

\item \textbf{Impedância e limite de linha não estão em dado aberto.} Estão, no
pacote \texttt{linha-transmissao} do ONS, com reatância em 100\% das linhas e
capacidade em MVA em 87\%. A afirmação incorreta aparecia em quatro pontos.

\item \textbf{Detectar painéis por satélite audita política de geração
distribuída.} Era a premissa original, e a seção~\ref{sec:fig2} a encerra.
\end{enumerate}
